Calcul de l'auteur : un terrain qui rapporte un loyer net R vaut R/r sans taxe et R/(r+t) avec une taxe annuelle au taux t de sa valeur ; r est le taux d'actualisation réel net de la croissance attendue des loyers.\par Lecture : les deux repères sont les taux qui, appliqués à la valeur des terrains bâtis hors administrations fin 2024, rapporteraient autant que la taxe foncière sur les propriétés bâties (0,7 \%) ou que celle-ci et les droits de mutation réunis (1 \%). Les prix actuels intègrent déjà ces impôts : une bascule à recette constante ne ferait pas baisser le prix moyen du terrain de l'ampleur lue sur les courbes, elle déplacerait la charge d'un bien à l'autre. La capitalisation complète est une hypothèse, que la section 4 confronte aux données.
